\chapter{Introduction}\label{ch:introduction}

\section{What akeyless-watchover is}
akeyless-watchover turns the \emph{governance} of an Akeyless account into code that is
reviewed, tested and delivered from one Git repository: one repository, one change, one
configuration delivered. It has two parts that work together:

\begin{description}[style=nextline,leftmargin=1.5em]
\item[The skill] An agent-neutral skill in the open Agent Skills format (\texttt{SKILL.md} plus
  reference documents), usable by any AI agent and model: natively by agents that support the
  format, through the bundled MCP server, or as a prompt pack (Chapter~\ref{doc-agents}). In a
  conversation, the agent translates an organisation diagram or a plain-language access request into Akeyless
  paths, roles, auth methods and associations, reviews changes, writes exact CLI commands and
  explains Akeyless authorization precisely. Chapter~\ref{ch:skill-commands} lists its commands.
\item[The kit] A tested repository template: a YAML model, generic Terraform that delivers it with
  the \texttt{akeyless-community/akeyless} provider, a one-time CI identity bootstrap, CI templates
  for GitHub and GitLab, and the \texttt{scripts/watchover} tool (offline tests, an access
  simulator, live verification, export of existing configuration to Terraform, Conjur policy
  translation, event-forwarder scaffolding and more). Chapter~\ref{ch:kit-commands} lists its commands.
\end{description}

\section{Scope}
In scope is the configuration \emph{inside} Akeyless: folders and paths, access roles and their
rules, auth methods and their bindings, management-plane permissions, account settings, custom
fields, key policies, event forwarders, PKI issuers and CA keys, and the Gateway's administrative
plane (allowed access, log forwarding, Conjur migration). Out of scope are deploying Gateways and
integrating consumers (applications, Kubernetes injectors, CI secret retrieval).

\section{Who this manual is for}
\begin{itemize}
\item \textbf{PAM and platform engineers} who own the Akeyless account: read everything, start
  with Chapter~\ref{ch:quickstart}.
\item \textbf{Business-unit engineers} who request or approve access: Chapter~\ref{ch:workflows}
  (recipes), the skill commands \texttt{change} and \texttt{verify}, and the model reference.
\item \textbf{Reviewers and auditors}: the RBAC model (Chapter~\ref{doc-rbac-model}), testing and
  verification, security considerations (Chapter~\ref{ch:security}) and the limitations.
\end{itemize}

\section{How this manual is organised}
\begin{description}[style=nextline,leftmargin=1.5em]
\item[Part~\ref{part:start}, Getting started] installation, using the skill with any AI agent, a guided first delivery, and recipes for everyday tasks.
\item[Part~\ref{part:concepts}, Concepts and design] the architecture, how Akeyless authorization
  works, and how to derive paths and roles from an organisation diagram.
\item[Part~\ref{part:reference}, Reference] every skill command and kit command with \emph{all}
  options and an example execution, every YAML key, the Terraform implementation, its variables and
  outputs, the CLI playbook, and how to choose between the REST API and the CLI.
\item[Part~\ref{part:operations}, Operations] pipeline, testing, brownfield import, Conjur
  migration, PKI, security, troubleshooting, maintenance and the known limitations.
\item[Appendices] exit codes, environment variables, repository layout, glossary and sources.
\end{description}

\section{Conventions}\label{sec:conventions}
\begin{itemize}
\item \texttt{monospace} is a command, path, option or YAML key. \texttt{<value>} is a placeholder;
  \texttt{[x]} is optional; \texttt{a|b} is a choice.
\item \textbf{VERIFY} marks a statement that no primary source confirms. All of them are collected
  in Chapter~\ref{doc-limitations}; check them on your account before relying on them.
\item Commands are run from the root of your governance repository unless a section says otherwise.
\end{itemize}

\subsection{How the example executions were produced}\label{sec:how-examples}
Every example in a grey terminal box was \emph{executed} by \texttt{docs/manual/build.py} when this
PDF was built, and the output is shown unedited except for long outputs, where the middle is cut
and marked \texttt{[... N lines omitted ...]}, and temporary paths, which are replaced by \texttt{.}.
The exit code is shown under each box.

\begin{note}[Offline demo back end]
No Akeyless account is used to build the manual. Kit commands that call Akeyless
(\texttt{verify}, \texttt{export}, \texttt{account-settings}, \texttt{token}) ran against a
stand-in: \texttt{docs/manual/demo/akeyless} plays the CLI and answers from the same fixtures the
test-suite uses, and a local HTTP endpoint plays the REST API. For \texttt{verify}, the stand-in
evaluates a copy of the model to which one manual console change was added, so the example shows
what a real drift finding looks like. On your account the commands are identical; only the data differs.
Where the command shown differs from the one executed (a production URL instead of the local
endpoint), the box says so.
\end{note}

The \emph{example sessions} of the skill commands (talking to your agent) are illustrative: the agent's
answers depend on the conversation. They show the shape of the output, abridged.
